\documentclass[12pt]{article}
\usepackage[utf8]{inputenc}
\usepackage[T1]{fontenc}
\usepackage{hyperref}
\title{Kayan Mahakam Language Resources for K-12 Curriculum Development}
\author{Compiled by Claude Sonnet 5.5 (claude-sonnet-5-5) for Andr\'as Kornai}
\date{}
\begin{document}
\maketitle

\section{Introduction}
Kayan Mahakam (ISO 639-3: xay; Glottocode: kaya1324) is an Austronesian language of the Kayanic branch, spoken in East Kalimantan, Indonesia. Wikipedia's page on the Kayan cluster cites about 35,000 speakers (1981--2007). That figure covers the cluster, not necessarily this variety alone. Very little was found, and nothing aimed at schools.

\section{Primary Educational Resources}
No primers, textbooks or school materials for Kayan Mahakam turned up. The language is not known to be taught in schools.

\section{Linguistic and Reference Materials}
\begin{itemize}
\item Kayan language (Borneo), Wikipedia: \url{https://en.wikipedia.org/wiki/Kayan_language_(Borneo)}
\item Mahakam languages, Wikipedia: \url{https://en.wikipedia.org/wiki/Mahakam_languages}
\item Languages of Kalimantan, Wikipedia: \url{https://en.wikipedia.org/wiki/Languages_of_Kalimantan}
\item Glottolog record: \url{https://glottolog.org/resource/languoid/id/kaya1324.bigmap.html}
\item Pacific Linguistics D-69 (the ``Holle Lists'' for Kalimantan, 1986). I did not confirm that it includes Kayan Mahakam entries: \url{https://openresearch-repository.anu.edu.au/bitstream/1885/144584/1/PL-D69.pdf}
\item 2012 conference paper noting that many languages of the area lack descriptions and data: \url{https://lingdy.aa-ken.jp/wp-content/uploads/2012/01/120217_soriente_inagaki_h2.pdf}
\item A 1911 comparative vocabulary of Kayan, Kenyah and Kalabit dialects, from Sarawak and not specific to Kayan Mahakam: \url{https://biodiversitylibrary.org/part/292444}
\end{itemize}

\section{Religious and Community Resources}
None were found in the searches.

\section{Digital and Audio Resources}
No audio, video or corpus was found.

\section{Curriculum Development Context}
No dedicated dictionary, grammar or curriculum material was found. Any curriculum would need new documentation and community consultation, supported by Indonesian-language schooling.
\end{document}